\section{Retaining the selected data through the final construction}
\label{sec:public-run-binding}

We finish the selection argument by specifying which common data are
retained. The issue is the order of the choices: the local estimates,
the numerical purchases and the limiting observables must all refer to
one family. The analytic hypotheses used here are those of
Sections~\ref{sec:common-funded-history},
\ref{sec:common-observable-purchases} and
\ref{sec:common-scale-selection}. This argument identifies their common
application; it supplies no replacement for an outstanding local estimate.

Fix a simple gauge group $G$ and a positive kernel penalty $b$.
Use the profile construction with support exponent exactly $100$ at
every stage, derivative order at least $6+j$ and spatial power at least
$2+j$ at stage $j$, as in Section~\ref{sec:ope-scale-purchases}. Denote the selected kernel by $K$, its stage profiles
by $\pi_j$, and its primary and secondary row-price families by $P$ and
$P'$. The flat-step, curvature-quadratic symmetry and dual-support estimates select their
floors from these data before the block lengths are chosen. Take their
maximum together with the original, comparison, treasury and period
floors, the marked-portfolio thresholds and the drift floors of
Section~\ref{sec:common-scale-selection}, and call it $n_j$. Raising the
floor to $n_j$ changes neither $K$ nor $\pi_j$ nor $P$. Select block
lengths $L_j\ge n_j$ and retain all these data when applying the
scheduling construction. The marked portfolio, with its sections and
regradings, was fixed before the profiles' price maps and the floors
(Section~\ref{sec:common-observable-purchases}); the run retains exactly
that portfolio, and no section or regrading is chosen after the
schedule.

Let $J(m)$ be the resulting stage clock, $\widehat\beta_m$ its hatted
coupling, and $H$ its quadratic horizon. Its actual block at step $m$
has length $L_{J(m+1)}$. For clarity, the two kinds of numerical payment
are different. If $U_{E,j}>0$ is the already selected width at fixed
index $E$, the countable purchase is
\[
 \forall E\ \exists m_E\ \forall m\ge m_E\ \forall j\le J(m+1):
 \quad U_{E,j}^{-1}\le 1+\log\log\widehat\beta_m.
\]
The same construction keeps $1\le1+\log\log\widehat\beta_m$ for
all $m$: the countable purchase includes $\exp(\exp E_j)$ with
$E_j\ge1$ in every $d_j$, so
$\widehat\beta_m\ge\beta_*\ge d_{j_*}\ge e^e$ and
$\log\log\widehat\beta_m\ge1$. After the lengths are fixed, let $A_j(L)$ and $B_j(L)$
be the two actual post-length costs and set
$F_j(L)=\max\{A_j(L),B_j(L)\}$. The all-stage purchase is
\begin{equation}\label{eq:public-native-allstage}
       F_{J(m+1)}(L)\le\widehat\beta_m\qquad(m\ge0).
\end{equation}
The countable version is retained as well. In particular, an eventual
statement for fixed $E$ is not being substituted for the all-stage
payment at $m=0$.

Here is the verification that the local suppliers apply to $H$ itself.
The selected levels have kernel $K$ and stage profiles $\pi_j$; the
horizon satisfies the row prices $P$. At step $m$, each component floor
is at most $n_{J(m+1)}\le L_{J(m+1)}$. The actual patch and the patch
used in computing $F_{J(m+1)}(L)$ agree, because a patch is determined
by its stage and block length. Thus \eqref{eq:public-native-allstage}
pays both actual costs. In the cost that has the form $e^{a_j}$, its
payment gives $a_j\le\log\widehat\beta_m$ by monotonicity of the
logarithm. The other payment supplies the flat-step delay. The entry delays of
the run are the final thresholds $d_j$ of
Section~\ref{sec:common-scale-selection}; they contain the terminal fee
threshold, the marked-portfolio delay, the startup and reset delays and
the loop-capacity thresholds listed there, so every schedule condition
of Sections~\ref{sec:ex-marked-portfolio-calibration},
\ref{sec:surviving-interval-shooting} and
\ref{sec:path-attachment-producer} that is stated as an entry delay
holds for this run's delays. The flat,
dual-support and curvature-quadratic symmetry suppliers therefore apply with the original
levels and $H$, at every allowed cutoff and level. No new horizon is
chosen in this application.

Now retain the finite-volume history, its campaign and the transition
witnesses obtained from this one construction. Denote them by
$R$, $T$ and $W$, respectively. Here $R$ names the full selected
history and scale data, $T$ names its campaign of parameter paths,
and $W$ keeps the transition and finite-volume witnesses for that
campaign. The family $C_R$ is the family of
Section~\ref{sec:original-history-portfolio-attachment}. In addition
to the fixed penalty $b$, keep the treasury and common-period estimates
and the order and power inequalities when making this simultaneous
choice. They are conclusions of the selection before the history is
packaged. One may retain all of them with the same existential witness;
forgetting them in a shorter statement does not prove that they can be
recovered from its remaining hypotheses.

There are also two distinct support-exponent facts. The stage
parameters used by the operator-product construction have exponent
exactly $100$. The larger budget used to bound local functions has
support exponent at least $100$. The first fact follows from the fixed
profile under reindexing by $J$; the second is a budget bound. Neither
requires changing the selected history or replacing its profile by an
increasing support-exponent profile.

The funded continuation is now evaluated at the history, reference
schedule, local sections, transition data and finite-volume witnesses
stored in $R$; its portfolio sections and regradings are the ones fixed
before the schedule, as above. Proposition~\ref{prop:ex-original-history-attachment}
identifies its output family with $C_R$. Its trace and coverage
conclusions consequently concern that same family. For the
operator-product estimates, keep $W$ as well as $R,T$: the finite
transition witnesses enter the service of
Section~\ref{sec:ope-history-service}. The fixed exponent, common
clock, halo and logarithmic estimates are those just selected.
Thus the common-history, observable and operator-product applications
are simultaneous applications of their canonical arguments.

\paragraph{The public family of this retained history.}
Write $\beta^R_{N,p,m}$ for the actual coupling of the stored history at
level $m\le N$, cutoff $N$ and period index $p$. Let $P_p$ be its selected
physical periods. To distinguish levels from stages, put
\[
 \ell_m=L_{J(m+1)},\qquad D_0=1,\qquad
 D_N=\prod_{m=0}^{N-1}\ell_m.
\]
Thus $\ell_m$ is the block used in the transition from level $m+1$ to
level $m$; it is not the stage length $L_m$ unless those indices agree.
All these data belong to $R$ and are fixed before an extraction order is
requested. Define its public Wilson family by retaining these periods,
using $P_pD_N$ sites on each axis, and setting
\begin{equation}\label{eq:public-family-coordinates}
 a_N=D_N^{-1}=e^{-s_N},\qquad
 \beta^{\rm pub}_{N,p}=\beta^R_{N,p,N}.
\end{equation}
The Haar measure and the Wilson action are those of
Section~\ref{ex-lattice}, in the fixed representation of $G$.

The insertion maps also retain the calibrated finite-volume data.
For a grade $d$ and cutoff-period pair $(N,p)$, let
$A^R_{d,N,p}$ be the terminal linear isomorphism of the coefficient space
stored in the transfer datum $S_R(d,N,p)$. Let
$\mathrm{Init}^R_{d,N,p}(w;x;U)$ be its initial local-section observable,
at site $x$ and configuration $U$, with coefficient $w$.
Here is its exact normalization. In the fixed grade-$d$ coefficient frame,
let $c_{d,\alpha}\in\mathbb N$ be the colour-factor count of coordinate
$\alpha$, and define the linear dilation
$(\Delta_d(t)w)_\alpha=t^{c_{d,\alpha}}w_\alpha$.
Write $\vartheta_N>0$ for the selected threshold and
$\mathcal L^R_d(w;x;U)$ for placement of the fixed local section with
coefficient $w$ at $x$. Then
\[
 \mathrm{Init}^R_{d,N,p}(w;x;U)
 =a_N^{-d}\mathcal L^R_d\bigl(\Delta_d(\vartheta_N^{-1})w;x;U\bigr).
\]
The dilation exponent is the coordinate's colour-factor count, whereas
the physical prefactor uses the grade $d$. These are the initial
observables used in Section~\ref{sec:selected-centered-service}.
The public token $(d,v)$ uses the same coefficient space and has insertion
\begin{equation}\label{eq:public-family-insertion}
 O^{\rm pub}_{N,p,(d,v)}(x;U)
 =\mathrm{Init}^R_{d,N,p}\bigl((A^R_{d,N,p})^{-1}v;x;U\bigr).
\end{equation}
In particular the terminal inverse is part of the insertion calibration;
the token $v$ is not substituted directly for its initial coefficient.
The formulas define the insertions at every cutoff using the stored
transfer datum there. Calibration and transport estimates are still
asserted only in their stated domains, in particular $N\ge t_R(d)$
when the reference-level hypothesis is required.

These identities follow from the retained construction, rather than
from a new choice of family. The schedule starts with denominator one
and multiplies it by its actual block $\ell_m$ at each level. Its scale
recurrence is $s_0=0$, $s_{m+1}=s_m+\log\ell_m$, so induction gives
$s_N=\log D_N$ and the first identity in
\eqref{eq:public-family-coordinates}. Since every $\ell_m\ge3$, the
positive integer denominators strictly increase. The stored history
starts at cutoff $N$ with bare Wilson density of coefficient
$\beta^R_{N,p,N}$ and zero vacuum parameter; this gives the second
identity. Finally, the finite insertion assigned by the transfer datum
is, by its definition, the initial local-section observable after its
terminal inverse. This is exactly \eqref{eq:public-family-insertion}.
Thus, for every $N,p$, every finite token tuple and every tuple of lattice
sites, the public finite moment is the moment of these same insertions
under the same normalized Wilson law of $C_R$. Neither a limiting
identification nor a replacement by the reference coupling
$\widehat\beta_N$ is used.

\paragraph{Conditional running consequence.}
Suppose the running estimate of
Section~\ref{sec:surviving-interval-shooting}, with the slope identified
as $\sigma=2b_0(G)>0$, holds for these actual coefficients: there is one
fixed intercept $\bar\beta$ such that for every $\eta>0$ there is
$m_0(\eta)$ with
\[
 \left|\beta^R_{N,p,m}-(\bar\beta+\sigma s_m)\right|\le\eta s_m
 \quad\text{for all }m\ge m_0(\eta),\ N\ge m,\ p\in\mathbb N.
\]
Apply the uniform extraction proposition of the running-coefficient
pilot to these coefficients, scale $s_m$ and intercept $\bar\beta$.
It gives a threshold independent of $N,p$ for their ratio to $s_m$.
Taking $m=N$ and using \eqref{eq:public-family-coordinates} yields
\[
 \forall\varepsilon>0\ \exists N_0\ \forall N\ge N_0\ \forall p:
 \quad
 \left|\frac{\beta^{\rm pub}_{N,p}}{\log(a_N^{-1})}-\sigma\right|
 <\varepsilon,
\]
which is \eqref{ex-running-input}. The denominator is positive on this
tail because $s_N\ge N\log3$. This is a specialization of the existing
uniform extraction argument, not a proof of its running-estimate input
or an identification with a coupling defined by a physical observable.

Finally this binding persists through the ordered extractions. Given
an initial extraction order $f$, first refine it to obtain the coupling
limits, then refine again to obtain the field limits. If the two
refinements are $g,w$, their combined order is
$f\circ g\circ w$, with the corresponding ordered period refinements.
Every coupling limit survives the second refinement, by the defining
limit property of a refinement. The phase, trace and operator-product
constructions are applied to the resulting single limiting table and
its single reconstructed system. Full-mesh convergence additionally
uses the full-mesh estimates on $C_R$; it is not inferred from the
existence of a mesh subsequence. These operations preserve the selected
history and its purchases while making the further choices required
by the limiting construction.
